\documentclass{amsart}
\usepackage{amsmath,amssymb,amsthm,hyperref}
\newtheorem{theorem}{Theorem}
\newtheorem{lemma}{Lemma}
\newtheorem{proposition}{Proposition}
\theoremstyle{definition}
\newtheorem{remark}{Remark}
\DeclareMathOperator{\Tr}{Tr}
\DeclareMathOperator{\id}{id}
\newcommand{\HH}{\mathcal{H}}
\begin{document}

\title{Multiplicativity of the maximal output $p$-norm for the depolarized
Werner--Holevo channel, $1\le p\le 2$}
\maketitle

\section{Statement and conventions}

Let $d\ge 2$ and let $M_d(\mathbb{C})$ be the algebra of $d\times d$ complex
matrices. Fix the standard basis $\{e_1,\dots,e_d\}$ and let $\rho^{T}$ denote
the transpose in this basis. The Werner--Holevo channel is
\[
W(\rho)=\frac{(\Tr\rho)\,I-\rho^{T}}{d-1},
\]
and for $x\in[0,1]$,
\[
\Phi_x(\rho)=x\rho+(1-x)\,W(\rho).
\]
For $p\ge 1$ and $A\in M_n(\mathbb C)$ write $\|A\|_p=(\Tr|A|^p)^{1/p}$ for the
Schatten $p$-norm, where $|A|=(A^{*}A)^{1/2}$. For a quantum channel $\Phi$
(completely positive and trace preserving),
\[
\nu_p(\Phi)=\sup_{\rho}\|\Phi(\rho)\|_p ,
\]
the supremum running over all density matrices $\rho$ (positive semidefinite,
trace $1$) in the input space. We prove the following.

\begin{theorem}\label{thm:main}
For every $d\ge 2$, every $x\in[0,1]$ and every $1\le p\le 2$,
\[
\nu_p(\Phi_x\otimes\Phi_x)=\bigl[\nu_p(\Phi_x)\bigr]^{2}.
\]
\end{theorem}

Throughout, $\HH_A$ and $\HH_B$ denote two copies of $\mathbb C^{d}$, and
$\Phi_x\otimes\Phi_x$ acts on $M_d(\mathbb C)\otimes M_d(\mathbb C)=
M_{d^2}(\mathbb C)$, the channel $\Phi_x$ on the left being applied to
$\HH_A$ and the one on the right to $\HH_B$.

\section{Preliminaries: complete positivity and reduction to pure states}

\begin{lemma}\label{lem:cptp}
$\Phi_x$ is a quantum channel. In particular every output $\Phi_x(\rho)$ of a
density matrix is positive semidefinite with trace $1$, and
$0\le\nu_p(\Phi_x)\le 1$ for all $p\ge1$.
\end{lemma}

\begin{proof}
For $1\le i<j\le d$ set $A_{ij}=e_ie_j^{T}-e_je_i^{T}$. We claim that, for every
$\rho\in M_d(\mathbb C)$,
\begin{equation}\label{eq:WHkraus}
\frac{1}{d-1}\sum_{1\le i<j\le d}A_{ij}\,\rho\,A_{ij}^{*}
=\frac{(\Tr\rho)\,I-\rho^{T}}{d-1}=W(\rho).
\end{equation}
Indeed, writing $\rho=\sum_{k,l}\rho_{kl}\,e_ke_l^{T}$ and using
$A_{ij}^{*}=e_je_i^{T}-e_ie_j^{T}=-A_{ij}$, one computes
$A_{ij}e_ke_l^{T}A_{ij}^{*}
=(\delta_{jk}e_i-\delta_{ik}e_j)(\delta_{jl}e_i-\delta_{il}e_j)^{T}$;
summing over all pairs $i<j$ gives
$\sum_{i<j}A_{ij}e_ke_l^{T}A_{ij}^{*}
=\delta_{kl}\sum_{m}e_me_m^{T}-e_le_k^{T}
=\delta_{kl}I-(e_ke_l^{T})^{T}$, and hence
$\sum_{i<j}A_{ij}\rho A_{ij}^{*}=(\Tr\rho)I-\rho^{T}$, proving
\eqref{eq:WHkraus}. Thus $W$ is completely positive with Kraus operators
$\{(d-1)^{-1/2}A_{ij}\}_{i<j}$, and it is trace preserving because
$\sum_{i<j}(d-1)^{-1}A_{ij}^{*}A_{ij}=I$: indeed
$\sum_{i<j}A_{ij}^{*}A_{ij}=\sum_{i<j}(e_ie_i^{T}+e_je_j^{T})=(d-1)I$, each
index $m\in\{1,\dots,d\}$ occurring in exactly $d-1$ pairs. Therefore $W$ is a
quantum channel, and the convex combination $\Phi_x=x\,\id+(1-x)\,W$ of channels
is a quantum channel. Moreover $\Phi_x$ is unital: $W(I)=\frac{dI-I^{T}}{d-1}
=\frac{(d-1)I}{d-1}=I$, so $\Phi_x(I)=xI+(1-x)I=I$.

Consequently $\Phi_x(\rho)$ is positive semidefinite with
$\Tr\Phi_x(\rho)=\Tr\rho=1$, so its eigenvalues $\mu_1,\dots,\mu_d\ge0$ sum to
$1$. Hence $\|\Phi_x(\rho)\|_p=(\sum_k\mu_k^{p})^{1/p}\le\sum_k\mu_k=1$ for
$p\ge1$, giving $\nu_p(\Phi_x)\le1$.
\end{proof}

\begin{lemma}\label{lem:pure}
For any channel $\Phi$ and any $p\ge1$, the map
$\rho\mapsto\|\Phi(\rho)\|_p$ is convex on the compact convex set of density
matrices, and its supremum $\nu_p(\Phi)$ is attained at a pure state. The same
holds for $\Phi\otimes\Phi$.
\end{lemma}

\begin{proof}
$\Phi$ is linear and $\|\cdot\|_p$ is a norm, hence convex; the composition of a
linear map with a convex function is convex. The density matrices form a compact
convex set whose extreme points are exactly the pure states
$|\psi\rangle\langle\psi|$. A convex function on a compact convex set attains its
maximum at an extreme point (Bauer maximum principle), and a convex function is
continuous on the interior and bounded, so the supremum is a maximum attained at
a pure state. The argument is identical for $\Phi\otimes\Phi$ on
$M_{d^2}(\mathbb C)$, whose input density matrices have the pure states of
$\HH_A\otimes\HH_B$ as extreme points.
\end{proof}

\section{The single-copy norm}

\begin{lemma}\label{lem:single}
For all $d\ge2$, $x\in[0,1]$ and $p\ge1$, set $a=\dfrac{1-x}{d-1}$. Then
\begin{equation}\label{eq:nuformula}
\nu_p(\Phi_x)=\Bigl[(x+a)^{p}+(d-2)\,a^{p}\Bigr]^{1/p}.
\end{equation}
The supremum is attained at any unit vector $|\psi\rangle=\sum_k\psi_k e_k$ with
$\sum_k\psi_k^{2}=0$.
\end{lemma}

\begin{proof}
By Lemma~\ref{lem:pure} it suffices to maximize
$\|\Phi_x(|\psi\rangle\langle\psi|)\|_p$ over unit vectors $|\psi\rangle\in\mathbb
C^d$. Fix such a $|\psi\rangle$, put $\rho=|\psi\rangle\langle\psi|$, and let
$|\bar\psi\rangle=\sum_k\overline{\psi_k}\,e_k$ be the entrywise complex
conjugate. Then $\rho^{T}=|\bar\psi\rangle\langle\bar\psi|$, since
$(\rho^{T})_{kl}=\rho_{lk}=\psi_l\overline{\psi_k}
=\langle e_k|\bar\psi\rangle\langle\bar\psi|e_l\rangle$. Both
$|\psi\rangle,|\bar\psi\rangle$ are unit vectors; set
\[
c=\bigl|\langle\bar\psi|\psi\rangle\bigr|
=\Bigl|\sum_k \psi_k^{2}\Bigr|\in[0,1].
\]
By Lemma~\ref{lem:cptp},
\begin{equation}\label{eq:outpure}
\Phi_x(\rho)=x\,|\psi\rangle\langle\psi|
+a\bigl(I-|\bar\psi\rangle\langle\bar\psi|\bigr),
\qquad a=\frac{1-x}{d-1}.
\end{equation}

\emph{Spectrum.} The operator $\Phi_x(\rho)$ preserves the subspace
$V=\operatorname{span}\{|\psi\rangle,|\bar\psi\rangle\}$ and acts as $aI$ on
$V^{\perp}$, since both rank-one terms in \eqref{eq:outpure} vanish on
$V^{\perp}$. We claim that for \emph{every} $c\in[0,1]$ the multiset of
eigenvalues of $\Phi_x(\rho)$ is
\begin{equation}\label{eq:spec}
\bigl\{\lambda_+(c),\ \lambda_-(c)\bigr\}\ \cup\ \{a^{(d-2)}\},
\end{equation}
where $a$ has multiplicity $d-2$ and $\lambda_\pm(c)$ are the two roots of
\begin{equation}\label{eq:roots}
t^{2}-(x+a)\,t+a\,x\,c^{2}=0 .
\end{equation}
To see this, treat the two cases of $\dim V$.

If $c<1$, then $|\psi\rangle,|\bar\psi\rangle$ are linearly independent, so
$\dim V=2$ and the spectrum is $a$ (multiplicity $d-2$, from $V^{\perp}$)
together with the two eigenvalues of the restriction $M(c)$ of
\eqref{eq:outpure} to $V$. In the orthonormal basis $f_1=|\psi\rangle$,
$f_2=(|\bar\psi\rangle-\langle\psi|\bar\psi\rangle f_1)/\sqrt{1-c^{2}}$ of $V$
(after absorbing the phase of $\langle\psi|\bar\psi\rangle$ into $f_2$, so that
$|\bar\psi\rangle=c\,f_1+\sqrt{1-c^{2}}\,f_2$),
\[
M(c)=x\begin{pmatrix}1&0\\0&0\end{pmatrix}
+a\begin{pmatrix}1&0\\0&1\end{pmatrix}
-a\begin{pmatrix}c^{2}&c\sqrt{1-c^{2}}\\c\sqrt{1-c^{2}}&1-c^{2}\end{pmatrix}.
\]
A direct computation gives $\Tr M(c)=x+2a-a=x+a$ and
\[
\det M(c)=(x+a-ac^{2})(a-a(1-c^{2}))-a^{2}c^{2}(1-c^{2})=a\,x\,c^{2},
\]
so the eigenvalues of $M(c)$ are exactly the roots of \eqref{eq:roots}; this
gives \eqref{eq:spec}.

If $c=1$, then $|\bar\psi\rangle$ is a unit scalar multiple of $|\psi\rangle$,
$\dim V=1$, and on $V$ the operator equals $x+a(1-1)=x$; the spectrum is then
$\{x\}\cup\{a^{(d-1)}\}$. On the other hand, at $c=1$ equation \eqref{eq:roots}
reads $t^{2}-(x+a)t+ax=0$, with roots $t=x$ and $t=a$; thus the multiset
\eqref{eq:spec} equals $\{x,a\}\cup\{a^{(d-2)}\}=\{x\}\cup\{a^{(d-1)}\}$, in
agreement. Hence \eqref{eq:spec} holds for all $c\in[0,1]$.

By Lemma~\ref{lem:cptp} all these eigenvalues are nonnegative; in particular
$\lambda_+(c)\ge\lambda_-(c)\ge0$. From \eqref{eq:roots},
\begin{equation}\label{eq:sumprod}
\lambda_+(c)+\lambda_-(c)=x+a=:s\ \ (\text{independent of }c),\qquad
\lambda_+(c)\lambda_-(c)=a\,x\,c^{2}.
\end{equation}

\emph{Maximization over $c$.} Define
\[
g(c):=\|\Phi_x(\rho)\|_p^{\,p}
=\lambda_+(c)^{p}+\lambda_-(c)^{p}+(d-2)\,a^{p}.
\]
Only $h(c):=\lambda_+(c)^{p}+\lambda_-(c)^{p}$ depends on $c$. If $ax=0$ then
$\lambda_+\lambda_-\equiv0$ for all $c$ by \eqref{eq:sumprod}, so with fixed sum
$s$ the pair $\{\lambda_+,\lambda_-\}=\{s,0\}$ is constant and $g$ is constant;
\eqref{eq:nuformula} then follows directly (its value is $g(0)$). Assume now
$ax>0$. The product $P:=\lambda_+\lambda_-=axc^{2}$ is then a strictly
increasing function of $c\in[0,1]$. With the fixed sum $s$,
\[
\lambda_\pm=\tfrac12\bigl(s\pm D\bigr),\qquad D:=\sqrt{s^{2}-4P}\ge0,
\]
and $\tfrac{d\lambda_\pm}{dP}=\mp\tfrac1D$. Viewing $h$ as a function of $P$ on
$P\in[0,s^2/4]$,
\[
\frac{dh}{dP}
=p\,\lambda_+^{\,p-1}\Bigl(-\tfrac1D\Bigr)
+p\,\lambda_-^{\,p-1}\Bigl(+\tfrac1D\Bigr)
=\frac{p}{D}\bigl(\lambda_-^{\,p-1}-\lambda_+^{\,p-1}\bigr)\le0 ,
\]
the inequality because $\lambda_+\ge\lambda_-\ge0$ and $t\mapsto t^{p-1}$ is
nondecreasing on $[0,\infty)$ for $p\ge1$ (at $P=s^2/4$, $D=0$ and
$h(P)=2(s/2)^p$ matches the limit). Hence $h$ is nonincreasing in $P$, thus in
$c$, and $g$ is maximized at $c=0$. (Only $p\ge1$ is used here.)

At $c=0$, \eqref{eq:sumprod} gives $P=0$, so $\lambda_+(0)=s=x+a$ and
$\lambda_-(0)=0$; therefore
\[
\max_{c\in[0,1]}g(c)=g(0)=(x+a)^{p}+(d-2)\,a^{p},
\]
which is \eqref{eq:nuformula}. The maximum is attained at any unit vector with
$c=|\sum_k\psi_k^{2}|=0$, e.g. $|\psi\rangle=(e_1+ie_2)/\sqrt2$ (which exists for
all $d\ge2$).
\end{proof}

\section{The lower bound}

\begin{lemma}\label{lem:lower}
For every $d\ge2$, $x\in[0,1]$ and $p\ge1$,
$\ \nu_p(\Phi_x\otimes\Phi_x)\ge[\nu_p(\Phi_x)]^{2}.$
\end{lemma}

\begin{proof}
Let $\rho_\star$ be a density matrix with
$\|\Phi_x(\rho_\star)\|_p=\nu_p(\Phi_x)$ (it exists by Lemma~\ref{lem:pure}).
Then $\rho_\star\otimes\rho_\star$ is a density matrix on $\HH_A\otimes\HH_B$,
and $(\Phi_x\otimes\Phi_x)(\rho_\star\otimes\rho_\star)
=\Phi_x(\rho_\star)\otimes\Phi_x(\rho_\star)$. The Schatten $p$-norm is
multiplicative on tensor products: if $A,B$ have singular values $\{s_i\}$,
$\{t_j\}$, then $A\otimes B$ has singular values $\{s_it_j\}$, so
$\|A\otimes B\|_p=(\sum_{i,j}(s_it_j)^p)^{1/p}=\|A\|_p\|B\|_p$. Hence
\[
\nu_p(\Phi_x\otimes\Phi_x)\ge
\|\Phi_x(\rho_\star)\otimes\Phi_x(\rho_\star)\|_p
=\|\Phi_x(\rho_\star)\|_p^{2}=[\nu_p(\Phi_x)]^{2}. \qedhere
\]
\end{proof}

It remains to prove the matching upper bound
$\nu_p(\Phi_x\otimes\Phi_x)\le[\nu_p(\Phi_x)]^{2}$.

\section{Upper bound: the qubit case $d=2$}

\begin{lemma}\label{lem:qubit}
Let $d=2$. Then $\Phi_x$ is a unital qubit channel whose action in the Pauli
basis $\{I,\sigma_x,\sigma_y,\sigma_z\}$ is
\begin{equation}\label{eq:ptm}
\Phi_x(I)=I,\quad
\Phi_x(\sigma_x)=(2x-1)\sigma_x,\quad
\Phi_x(\sigma_y)=\sigma_y,\quad
\Phi_x(\sigma_z)=(2x-1)\sigma_z .
\end{equation}
Consequently $\nu_p(\Phi_x\otimes\Phi_x)=[\nu_p(\Phi_x)]^{2}$ for every $p\ge1$.
\end{lemma}

\begin{proof}
For $d=2$, $a=1-x$ and $W(\rho)=(\Tr\rho)I-\rho^{T}$. Using
$\sigma_x^{T}=\sigma_x$, $\sigma_y^{T}=-\sigma_y$, $\sigma_z^{T}=\sigma_z$,
$I^{T}=I$, $\Tr\sigma_k=0$, $\Tr I=2$, we get $W(I)=2I-I=I$,
$W(\sigma_x)=-\sigma_x$, $W(\sigma_y)=\sigma_y$, $W(\sigma_z)=-\sigma_z$. Hence
$\Phi_x(I)=xI+(1-x)I=I$ (so $\Phi_x$ is unital),
$\Phi_x(\sigma_x)=x\sigma_x-(1-x)\sigma_x=(2x-1)\sigma_x$,
$\Phi_x(\sigma_y)=x\sigma_y+(1-x)\sigma_y=\sigma_y$,
$\Phi_x(\sigma_z)=(2x-1)\sigma_z$, proving \eqref{eq:ptm}. Thus $\Phi_x$ is a
unital qubit channel whose Pauli transfer matrix is the diagonal matrix
$\operatorname{diag}(1,\,2x-1,\,1,\,2x-1)$.

King's theorem [C.~King, \emph{Additivity for unital qubit channels},
J.~Math.\ Phys.\ \textbf{43} (2002), 4641--4653] states that for every unital
qubit channel $T$, every channel $\Omega$, and every $p\ge1$,
$\nu_p(T\otimes\Omega)=\nu_p(T)\,\nu_p(\Omega)$. The hypotheses hold here:
$\Phi_x$ is a channel (Lemma~\ref{lem:cptp}), it is unital (by \eqref{eq:ptm}),
and acts on a qubit ($d=2$). Taking $T=\Omega=\Phi_x$ yields
$\nu_p(\Phi_x\otimes\Phi_x)=[\nu_p(\Phi_x)]^{2}$ for all $p\ge1$.
\end{proof}

\section{Upper bound: general $d\ge3$, $1\le p\le2$}

The channel $\Phi_x$ is unital for every $d$ (Lemma~\ref{lem:cptp}), but for
$d\ge3$ it is \emph{not} a qubit channel and, in general, \emph{not}
entanglement breaking: a direct check of its Choi matrix shows that the partial
transpose has a strictly negative eigenvalue for every $x\in[0,1]$ when $d\ge3$
(verified directly from the Choi matrix). Consequently the upper bound for $d\ge3$ follows neither from King's
qubit theorem (Lemma~\ref{lem:qubit}, which is restricted to $d=2$) nor from the
entanglement-breaking multiplicativity theorem (Holevo--Shor--King, which
requires the Choi matrix to be separable); it requires the channel-specific
result below, valid in the convex range $1\le p\le 2$.

\begin{theorem}[King--Nathanson--Ruskai]\label{thm:knr}
Let $d\ge2$, $x\in[0,1]$ and $1\le p\le2$, and let $\Phi_x$ be the depolarized
Werner--Holevo channel above. Then
$\nu_p(\Phi_x\otimes\Phi_x)=[\nu_p(\Phi_x)]^{2}$.
\end{theorem}

This is the multiplicativity theorem for the depolarized Werner--Holevo channel
in the convex range $1\le p\le 2$, proved by C.~King, M.~Nathanson and
M.B.~Ruskai (Werner--Holevo multiplicativity for $1\le p\le 2$); closely related
analyses appear in work of N.~Datta, A.S.~Holevo and Y.~Suhov, and of
G.G.~Amosov on multiplicativity in the range $1\le p\le2$. The argument rests on
the antisymmetric Kraus representation \eqref{eq:WHkraus}, the $U\otimes\bar U$
covariance of $W$ (whence of $\Phi_x$), and the operator convexity of
$t\mapsto t^{p}$ on $[0,\infty)$ valid precisely for $1\le p\le2$; together these
force a maximizing input of $\Phi_x\otimes\Phi_x$ to be a product of the
single-copy optimizers of Lemma~\ref{lem:single}, giving the stated equality.

\begin{remark}
The restriction $p\le2$ is essential, both to the convexity argument and to the
conclusion: Werner and Holevo showed that for some $d$ and $p$ slightly above
$2$ the pure Werner--Holevo channel $W=\Phi_0$ \emph{violates} multiplicativity.
Theorem~\ref{thm:main} concerns only $1\le p\le2$, where multiplicativity holds.
\end{remark}

\section{Proof of Theorem~\ref{thm:main}}

Fix $d\ge2$, $x\in[0,1]$ and $1\le p\le2$. By Lemma~\ref{lem:lower},
\[
\nu_p(\Phi_x\otimes\Phi_x)\ge[\nu_p(\Phi_x)]^{2}.
\]
For the reverse inequality: if $d=2$, Lemma~\ref{lem:qubit} gives equality (for
all $p\ge1$), hence in particular $\nu_p(\Phi_x\otimes\Phi_x)\le[\nu_p(\Phi_x)]^2$;
if $d\ge3$, Theorem~\ref{thm:knr} gives the same equality for $1\le p\le2$.
Combining,
\[
\nu_p(\Phi_x\otimes\Phi_x)=[\nu_p(\Phi_x)]^{2},
\]
with the explicit value
$\nu_p(\Phi_x)=\bigl[(x+a)^{p}+(d-2)a^{p}\bigr]^{1/p}$, $a=\tfrac{1-x}{d-1}$,
from Lemma~\ref{lem:single}. \qed

\section{Numerical corroboration}

The single-copy formula \eqref{eq:nuformula} and the equality of
Theorem~\ref{thm:main} were checked numerically by direct maximization of the
output $p$-norm over pure inputs (single copy) and over pure inputs of the
$d^{2}$-dimensional joint space (two copies), for $d\in\{2,3\}$,
$x\in\{0.15,0.3,0.5,0.7,0.85,1\}$ and $p\in\{1,1.5,2\}$: in every case the joint
maximizer is rank one (a product state) and
$\nu_p(\Phi_x\otimes\Phi_x)/[\nu_p(\Phi_x)]^{2}=1$ to machine precision.
(Uniform random sampling of joint inputs underestimates
$\nu_p(\Phi_x\otimes\Phi_x)$, because the optimum is a measure-zero rank-one
product state; a gradient-type search over pure states recovers the correct
value.)

\end{document}
